\documentclass[main.tex]{subfiles}


\begin{document}

%from pagenumber to up
% \phantomsection 
% \hypertarget{toc}{}

 

 

% номер секции вручную
\setcounter{section}{43
}
\addtocounter{section}{-1} 

\section{Свободные колебания в контуре}


Свободные (\emph{собственные}) электромагнитные колебания, происходящие в колебательном контуре, оказываются \emph{гармоническими}. Так, \textbf{заряд конденсатора} меняется во времени по закону \emph{косинуса} (или \emph{синуса}):
\begin{equation}
    q=q_\text{max}\cos(\omega t + \varphi),
\end{equation}
где $q_\text{max}$~--- \emph{амплитуда заряда} (наибольшее значение заряда), $\omega$~--- циклическая частота (или круговая частота),
% \footnote{Свободные колебания в контуре совершаются с так называемой собственной циклической частотой~$\omega_\text{с}$ (или с линейной частотой~$\nu_\text{с}$).}
$\varphi$~--- начальная фаза.

\textbf{Ток катушки} и \textbf{напряжение конденсатора} можно вычислять через заряд конденсатора:
\begin{equation}
    I=q',\quad U=\frac{q}{C},
\end{equation}
где $q'$~--- производная заряда по времени, $C$~--- емкость конденсатора.

На рис.\,\ref{fig:p167} изображены графики колебаний заряда конденсатора и тока катушки в контуре (в одних координатных осях).

\begin{figure}[H]
    \centering

    \begin{tikzpicture}[font=\small,            line cap=round,                             >={Stealth[inset=0pt,round,length=3mm, angle'=20]},vec/.style={>={Stealth[inset=0pt,round,length=3.5mm, angle'=20]}}]


\path (.9,3) -- (.3,3)
foreach \y in {.2,.6,...,1.2}
{ -- (.9,{3-\y}) -- (.3,{3-\y-.2}) }
-- (.6,{3-1.2-.2/2}) coordinate (cir)
;
\path (cir);
\pgfgetlastxy{\cx}{\cy}

\draw[->] (0,{\cy-.3cm}) ++ (0,-2) --++ (0,4) node[left,black] {$q,I$};
\draw[->] (0,{\cy-.3cm}) --++ (11,0) node[below,black] {$t$};
\node[below left] (O) at (0,{\cy-.3cm}) {$O$};

\def\xmax{(3-.9-.1/2-(3-1.5-.2/2))}
\def\er{6.2*pi}
\draw[NavyBlue,thick,domain=0:\er,yshift={\cy-.3 cm},xscale=.5,samples=1000] plot (\x,{\xmax*cos(\x r)}) node [right,black] {$q$};
\draw[red,thick,domain=0:\er,yshift={\cy-.3 cm},xscale=.5,samples=1000] plot (\x,{-1.5*sin(\x r)}) node [right,black] {$I$};
% \draw[red,dashed,thick,domain=pi:4*pi,yshift={\cy-.3 cm},,xscale=.5,samples=1000] plot (\x,{\xmax*cos(\x r)});
\path[] (pi,{\cy-.3 cm+\xmax*1cm}) --++ (0,{-\xmax}) node[circle,fill=black,inner sep=0pt,minimum size=1mm,label={below left,inner xsep=0}:$T$] {};
\path[] (2*pi,{\cy-.3 cm+\xmax*1cm}) --++ (0,{-\xmax}) node[circle,fill=black,inner sep=0pt,minimum size=1mm,label={below left,inner xsep=0}:$2T$] {};



% %dash
% \draw[gray,densely dashed] (.6,{3-.9-.1/2}) --++ (5.4,0) node[above left,black] {$x_\text{max}$};
% \draw[gray,densely dashed] (2.1,{3-1.5-.2/2}) --++ (3.9,0);
% \draw[gray,densely dashed] (3.6,{3-2.1-.3/2}) --++ (2.4,0) node[below left,black] {$-x_\text{max}$};


    \end{tikzpicture}
\caption{Графики колебаний заряда и тока в контуре}
\label{fig:p167}
\end{figure}

Ток и заряд колеблются с периодом~$T$.
% (или с так называемой собственной частотой контура~$\nu_\text{с}$). 
При этом нули заряда приходятся на максимумы или минимумы тока; и наоборот, нули тока соответствуют максимумам или минимумам заряда.

\textbf{Период колебаний} в контуре находят так (\emph{формула Томсона}):
\begin{equation}
    T=2\pi\sqrt{LC},
\end{equation}
где $L$ и $C$~--- индуктивность катушки и емкость конденсатора. 

При решении задач удобно использовать понятие полной энергии контура.

\textbf{Полная энергия контура} ($W$ [Дж])~--- это сумма энергий катушки и конденсатора:
\begin{equation}
    W=W_L+W_C.
\end{equation}
%
% Так как предполагается, что сопротивление контура равно нулю (в общем энергия )
%
Следующий закон обычно применяется при решении задач, сформулированных в терминах энергии.
%
\begin{law}
    \textbf{Закон сохранения полной энергии контура.} Если сопротивление контура равно нулю (энергия перераспределяется только между катушкой и конденсатором), то полная энергия контура сохраняется:
    \begin{equation}
        W_1=W_2=\ldots,
    \end{equation}
    где $W_1$ и $W_2$~--- полные энергии контура в первом и втором состояниях. 
\end{law}




 
\end{document}